\scanpage{13}\mspage{10}
\begingroup
\makeatletter\tagsleft@true\makeatother
8. \uline{Algorithm} \uline{continued}, \uline{the} \uline{last five} \uline{cases}. Assume now that $\boxed{\mathfrak p^3\mid b_6}$. Then we can change coördinates so that $\boxed{\mathfrak p\mid a_1\text{ and }a_2,\quad \mathfrak p^2\mid a_3\text{ and }a_4,\quad \text{and }\mathfrak p^3\mid a_6}$. This being done, consider the polynomial
\[
 P(T)=T^3+a_{2,1}T^2+a_{4,2}T+a_{6,3},
 \tag*{(8.1)}
\]
where the notation is as explained in the proof of step \textcircled{4}, page $9$. Our equation now becomes
\[
 \pi y_2^2+\pi a_{1,1}x_1y_2+\pi a_{3,2}y_2=P(x_1),
 \tag*{(8.2)}
\]
and our algorithm has three branches, according to the multiplicities of the roots of the congruence $P(T)\equiv0$:

\reversemarginpar
\margin{($1^{\mathrm{st}}$ branch)}\textcircled{6} If $P(T)\equiv0$ has three distinct roots, then we have Type $\mathrm I_0^{*}$, and
\[
 f=v(\Delta)-4,\quad\text{and}\quad c=1+(\text{number of roots of }P(T)\equiv0\text{ in }k)
\]

\margin{($2^{\mathrm d}$ branch)}\textcircled{7} If $P(T)\equiv0$ has one simple and one double root, then Type $\mathrm I_{\nu}^{*}$, $\nu>0$, and
\[
 f=v(\Delta)-4-\nu,\qquad c=2\text{ or }4,
\]
where $\nu$ and $c$ are obtained by the following procedure:

\margin{\rotatebox{90}{subprocedure branch \textcircled{7}}%
\figure{\begin{tikzpicture}[x=1cm,y=1cm]
\draw[line width=.8pt,line join=round] (1.25,0) -- (.22,0) -- (.22,-6.2);
\draw[line width=.8pt,line join=round] (.12,-5.94) -- (.22,-6.2) -- (.40,-5.94);
\end{tikzpicture}}}%
Translate $x$, so that $T\equiv0$ is a double root of $P(T)\equiv0$. Then $\mathfrak p^3\mid a_4$, $\mathfrak p^4\mid a_6$, $\mathfrak p^2\nmid a_2$ and (8.2) becomes
\[
 y_2^2+\pi a_{1,1}x_2y_2+a_{3,2}y_2
 =\pi^2x_2^3+\pi a_{2,1}x_2^2+\pi a_{4,3}x_2+a_{6,4}.
 \tag*{\fbox{7.1}}
\]

If $Y^2+a_{3,2}Y-a_{6,4}\equiv0$ has distinct roots, then
\[
 \nu=1,\quad\text{and}\quad c=\begin{cases}4&\text{if roots in }k.\\2&\text{otherwise }.\end{cases}
\]

If $Y^2+a_{3,2}Y-a_{6,2}\equiv0$ has a double root, we can translate $y$ so that the root is $Y\equiv0$. Then $\mathfrak p^3\mid a_3$, $\mathfrak p^5\mid a_6$, and our equation becomes
\[
 \pi y_3^2+\pi a_{1,1}x_2y_3+a_{3,3}y_3
 =\pi x_2^3+a_{2,1}x_2^2+a_{4,3}x_2+a_{6,5}
 \tag*{\fbox{7.2}}
\]

If $a_{2,1}X^2+a_{4,3}X+a_{6,5}\equiv0$ has distinct roots, then
\[
 \nu=2,\ \text{and}\ c=\begin{cases}4&\text{if roots in }k\\2&\text{otherwise}.\end{cases}
\]
\normalmarginpar
\hfill(continued)
\endgroup
% archive sheet number: 12
